\section{Audit and Proxy-Corrected Estimation}
\label{sec:v8-audit-certificate}

\subsection{Audit Units}

The audit treats the realized tree as a finite population. Every
node in \(T\) whose C1/C2/C3 losses define the local-law channel is
a unit in the population, and the audit samples some of them under
a known design. We index audit units by \(i\) (the
survey-statistics convention; each \(i\) is a node \(v\) from
Section~\ref{sec:v8-ctree-math}). Unit \(i\) carries a represented
raw span \(S(i)=x_v\), a produced state \(s_i=z_v\), a law channel
(C1, C2, or C3), and a logged marginal inclusion probability
\(\pi_i>0\) recording how often the audit design samples this node.

Two diagnostics come out of the audit. The first is the violation
rate per law channel: how often C1 fails, how often C2 fails, how
often C3 fails. This tells the analyst which kind of failure is
driving the residual: leaves dropping evidence, states changing
when summarized again, or merges losing the interaction between
siblings. The second is the mean distortion across nodes, which is
what feeds the certificate. A run can have many small violations
whose root impact is limited, or a few large ones that move the
score a lot. Reporting both quantities lets the analyst tell the
two cases apart.

\subsection{Dense Proxy, Sparse Oracle}

We cannot compute \(L_{\mathrm{law}}^\ast\) in
Equation~\eqref{eq:v8-local-law-estimand} directly, because it
requires an \(f^\ast\) call at every node and oracle calls are
scarce. Long-document settings typically expose oracle labels at
the root (e.g., expert-survey means on the whole document) and
leave internal nodes without a direct oracle. What is available at
internal nodes is some accessible substitute: a cached teacher
trace, an aggregated label from a span-grounded coding scheme, or
the proxy \(\widehat f\). The fix is to use the proxy densely and
correct it with sampled oracle calls.

The dense local quantity is the same C1/C2/C3 comparison measured
through the proxy:
\[
  \widehat\ell_v
  =
  \ell_{\mathrm{law}}(v;\widehat f,g).
\]
\(\widehat\ell_v\) gives a training signal at every node, but the
proxy is not the oracle. The gap between them at node \(v\) is the
node bias:
\[
  b_v=\widehat\ell_v-\ell_v^\ast .
\]
Sampled \(f^\ast\) calls estimate \(b_v\) on the nodes they cover, and
inverse-propensity weighting carries that estimate to the population.

Let \(R_v\in\{0,1\}\) indicate whether node \(v\) is sampled and
evaluated with \(f^\ast\), and let \(\pi_v>0\) be the logged marginal
probability that the audit design samples \(v\). The corrected node
loss is
\begin{equation}
\label{eq:v8-corrected-node-loss}
  \ell_v^{\mathrm{corr}}
  =
  \widehat\ell_v
  +
  \frac{R_v}{\pi_v}
  \left(\ell_v^\ast-\widehat\ell_v\right).
\end{equation}
Unsampled nodes contribute proxy loss. Sampled nodes add back the
observed oracle--proxy difference, scaled up by \(1/\pi_v\) so the
correction stands in for the population mean. This is the standard
inverse-propensity correction: rare nodes are sampled less, so the
few times they are sampled they count more.

\begin{proposition}[Corrected Node Loss Is Unbiased for the Oracle
Loss]\label{prop:v10-corrected-loss-unbiased}
For each node \(v\) with logged inclusion probability \(\pi_v>0\),
\[
  \mathbb E\!\left[\ell_v^{\mathrm{corr}}\mid T\right]
  =\ell_v^\ast .
\]
\end{proposition}

\begin{proof}[Proof sketch]
Assume the audit design samples node \(v\) with logged inclusion
probability \(\pi_v>0\). Conditional on the realized tree, the
sampling indicator \(R_v\) is Bernoulli with mean \(\pi_v\), so the
correction factor \(R_v/\pi_v\) has mean \(1\). Take expectations of
the corrected node loss in
Equation~\eqref{eq:v8-corrected-node-loss}: the proxy term
\(\widehat\ell_v\) is non-random and stays put; the correction term
\((R_v/\pi_v)(\ell_v^\ast-\widehat\ell_v)\) has mean
\(\ell_v^\ast-\widehat\ell_v\); summing the two leaves
\(\ell_v^\ast\). Strict positivity of \(\pi_v\) is the only place the
argument uses; the design can otherwise be adversarial. The full
argument and the variance bound are in
Appendix~\ref{app:v8-objective-ipw}.
\end{proof}

The corrected local-law channel sums the corrected node losses with
the same depth weights as the oracle channel:
\[
  L_{\mathrm{law}}^{\mathrm{corr}}(g,T)
  =
  \sum_{v\in V(T)}\gamma^{d(v)}\ell_v^{\mathrm{corr}} .
\]
Plugging \(L_{\mathrm{law}}^{\mathrm{corr}}\) into the
theorem-facing objective from Section~\ref{sec:v8-objective} gives
the trainable version of the C-TreePO objective:
\begin{equation}
\label{eq:v8-corrected-objective}
  J_\lambda^{\mathrm{corr}}(g;T)
  =
  (1-\lambda)L_{\mathrm{root}}(g)
  +
  \lambda L_{\mathrm{law}}^{\mathrm{corr}}(g,T).
\end{equation}
\(J_\lambda^{\mathrm{corr}}\) estimates the oracle objective
\(J_\lambda^\ast\) from
Equation~\eqref{eq:v8-main-objective}. With no audit calls
(\(R_v=0\) at every node), the corrected channel reduces to the
dense-proxy channel; sampled \(f^\ast\) calls add the correction.

\subsection{HT Distortion Certificate}

The corrected loss gives a trainable objective. We also need a
finite-sample bound on the downstream drift \(G_{\mathrm{meth}}\),
built from the same audit. Let \(\mathcal U\) be the finite
population of sample-eligible tree units, \(N=|\mathcal U|\),
and let
\[
  D_i^J
  =
  d_Y\!\left(J(s_i),J(S(i))\right)
\]
be the distortion measured by an accessible scorer \(J\). \(J\) is
whichever scorer the audit instantiates: in the corrected-loss form
above it is the proxy \(\widehat f\), and on a sampled subset of nodes
\(J\) can be the oracle \(f^\ast\) itself. If \(Z_i\) is the audit
inclusion indicator and \(\pi_i\) is the logged inclusion probability,
the Horvitz--Thompson estimator is
\begin{equation}
\label{eq:v8-ht-distortion}
  \widehat{\Delta}_{\mathrm{HT}}^{J}
  =
  \frac{1}{N}
  \sum_{i\in\mathcal U}
  \frac{Z_i}{\pi_i}D_i^J .
\end{equation}

\begin{proposition}[HT Distortion Estimator Is Unbiased]
\label{prop:v10-ht-unbiased}
Assume strict positivity (\(\pi_i>0\) for every \(i\in\mathcal U\)) and
that \(\pi_i\) is the correct logged marginal inclusion probability
for unit \(i\). Then
\[
  \mathbb E\!\left[\widehat{\Delta}_{\mathrm{HT}}^{J}\mid T\right]
  =
  \frac{1}{N}\sum_{i\in\mathcal U}D_i^J ,
\]
the finite-population mean of the per-unit distortions. Unequal
propensities affect variance and effective sample size, not the
target being estimated.
\end{proposition}

\begin{proof}[Proof sketch]
Assume strict positivity (\(\pi_i>0\) for every unit) and that
\(\pi_i\) is the correct logged marginal inclusion probability of
unit \(i\). The argument is the standard Horvitz--Thompson identity
applied unit by unit. For unit \(i\), the inclusion indicator
\(Z_i\) has mean \(\pi_i\), so \(Z_i/\pi_i\) has mean \(1\) and
\((Z_i/\pi_i)D_i^J\) has mean \(D_i^J\). Summing across the finite
population and dividing by \(N\) gives the claimed finite-population
mean. Strict positivity is what makes \(Z_i/\pi_i\) well defined and
prevents any unit from being unreachable. Unequal propensities
change the variance through the same identity but leave the mean
alone, since the per-unit factors have not changed. The IPW stress
check in Section~\ref{sec:v10-ipw-check} reports empirical coverage
on adversarial finite populations. The full argument with the
matching variance bound is in Appendix~\ref{app:v8-objective-ipw}.
\end{proof}

The certificate has four terms. The first is the audit's estimate of
the representation distortion, transported through the downstream
method. The other three are correction terms for calibration,
sampling error, and clipping.

\begin{proposition}[Finite-Sample Certificate]
\label{prop:v10-finite-sample-certificate}
Under strict positivity and correct logged propensities for the
audit design, oracle-compatibility of the downstream method with
transport constant \(C_{\mathrm{meth}}\), bounded clipping radius,
and declared calibration, sampling, and clipping envelopes on an
event \(\mathcal E\), the downstream drift admits the four-term
decomposition
\begin{equation}
\label{eq:v8-finite-sample-certificate}
  |G_{\mathrm{meth}}|
  \le
  C_{\mathrm{meth}}\widehat{\Delta}_{\mathrm{HT}}
  +
  B_{\mathrm{cal}}
  +
  B_{\mathrm{est}}
  +
  B_{\mathrm{clip}} .
\end{equation}
If the component envelopes hold with total failure probability at
most \(\delta\), then the certificate holds with probability at least
\(1-\delta\).
\end{proposition}

\begin{proof}[Proof sketch]
Assume strict positivity and correct logged propensities (so that
Proposition~\ref{prop:v10-ht-unbiased} applies),
oracle-compatibility of the downstream method through the transport
constant \(C_{\mathrm{meth}}\), and a bounded clipping radius on
the inverse-propensity weights. Also assume that the calibration,
sampling, and clipping envelopes hold on the event \(\mathcal E\).
The argument adds and subtracts two intermediate quantities inside
the transport bound: the accessible-judge population distortion and
the unclipped HT estimate. The four resulting pieces are the
method-transported audit estimate
\(C_{\mathrm{meth}}\widehat\Delta_{\mathrm{HT}}\); the calibration
gap \(B_{\mathrm{cal}}\) between the proxy scorer and the oracle;
the sampling deviation \(B_{\mathrm{est}}\) of the HT estimator from
its mean; and the clipping bias \(B_{\mathrm{clip}}\) introduced
when large inverse-propensity weights are truncated for variance
control. The deterministic inequality holds on \(\mathcal E\). The
high-probability version assigns failure budgets to the component
envelopes and combines them by a union bound. Full derivation in
Appendix~\ref{app:v8-full-proofs}, with the per-component bounds
spelled out in Appendix~\ref{app:v8-objective-ipw}.
\end{proof}

\(C_{\mathrm{meth}}\widehat{\Delta}_{\mathrm{HT}}\) is the sampled
distortion times the method transport constant from
Section~\ref{sec:v8-objective}. \(B_{\mathrm{cal}}\) is the gap
between the accessible scorer and \(f^\ast\): the part of error that
more sampled labels cannot remove. \(B_{\mathrm{est}}\) is sampling
uncertainty after method transport. \(B_{\mathrm{clip}}\) is bias
introduced by clipping inverse-propensity weights or residuals.

The Benoit replication below reports the root-observed tier and
teacher-trace local diagnostics. A completed human node audit would
instantiate \(D_i^J\) with span-level expert or quasi-sentence
aggregates, log the node propensities, and report
Equation~\eqref{eq:v8-finite-sample-certificate} as the realized
certificate. Section~\ref{sec:v10-controlled-checks} supplies the
separate estimator stress check for the HT/IPW component.
